%%%%%%%%%%%%%%%%%%%%%%%%%%%%%%%%%%%%%%%%%%%%%%%%%%%%%%%%%%%%%%%%%%%%%%%%%%%%%%%%
%%
\cleardoubleoddpage%  Make sure to start each chapter on a new odd page
\chapter{Results and Analysis}
\label{sec:results}

This chapter presents the experimental results from evaluating \ac{xLSTM} variants in the \ac{BIOT} architecture for \ac{EEG} event classification using the \ac{TUEV} dataset. The experimental design systematically compares four architectural variants (\ac{sLSTM}, \ac{mLSTM}, combined \ac{mLSTM}+\ac{sLSTM}, and a Linear Transformer baseline) across three temporal window configurations (5\,s, 7\,s, 9\,s). Each configuration was run five times with different random seeds (10, 42, 123, 456, 999) to ensure statistical robustness.

\subsection{Model Selection Criterion}

\textbf{All results use checkpoints selected by balanced accuracy, not validation loss.} The \ac{TUEV} dataset exhibits severe class imbalance (66.77\% majority class), causing validation loss to plateau and fluctuate after initial convergence. Models can minimize cross-entropy loss by biasing toward the majority class while achieving poor balanced accuracy. Since balanced accuracy equally weights all classes, it is the appropriate selection criterion for this imbalanced multi-class task. All reported metrics, confusion matrices, and statistical comparisons use balanced-accuracy checkpoints exclusively.

\section{Overall Performance Comparison}

We analyze four key metrics (see Section~\ref{sec:evaluation_metrics} for definitions): balanced accuracy (robust to class imbalance), Cohen’s~\(\kappa\) (agreement beyond chance), weighted F1 (precision–recall harmonic mean weighted by support), and \ac{AUCPR} Macro (ranking quality for rare events). Overall accuracy is not emphasized because it is misleading under imbalance; when shown, it serves only for completeness in statistical tests.

\subsection{Complete Performance Results}

Table~\ref{tab:metrics_comparison} summarizes performance across all 60 runs (4 architectures~\(\times\) 3 windows~\(\times\) 5 seeds). Values are mean~\(\pm\)~standard deviation over the five seeds; bold indicates the best score within a given window length.

\begin{table}[htbp]
\centering
\caption{Performance comparison across architectures and temporal window configurations. Values are mean~\(\pm\)~std over five seeds using balanced-accuracy checkpointing. Boldface indicates the best value per window.}
\label{tab:metrics_comparison}
\footnotesize
\setlength{\tabcolsep}{3pt}
\begin{tabular}{l | c | c | c | c | c}
\hline
	extbf{Architecture} & \textbf{Win.} & \textbf{Bal. Acc.} & \textbf{Cohen \(\kappa\)} & \textbf{F1 (W)} & \textbf{\ac{AUCPR} (Macro)} \\
\hline
Linear Transf. & 5s & $0.521 \pm {\scriptsize 0.041}$ & $0.457 \pm {\scriptsize 0.049}$ & $0.712 \pm {\scriptsize 0.034}$ & $0.468 \pm {\scriptsize 0.013}$ \\
Linear Transf. & 7s & $0.501 \pm {\scriptsize 0.043}$ & $0.465 \pm {\scriptsize 0.030}$ & $0.717 \pm {\scriptsize 0.022}$ & $0.465 \pm {\scriptsize 0.039}$ \\
Linear Transf. & 9s & $0.485 \pm {\scriptsize 0.037}$ & $0.425 \pm {\scriptsize 0.048}$ & $0.691 \pm {\scriptsize 0.031}$ & $0.453 \pm {\scriptsize 0.026}$ \\
\hline
\ac{mLSTM} & 5s & $0.509 \pm {\scriptsize 0.031}$ & $0.439 \pm {\scriptsize 0.037}$ & $0.707 \pm {\scriptsize 0.017}$ & $0.487 \pm {\scriptsize 0.030}$ \\
\ac{mLSTM} & 7s & $\mathbf{0.507 \pm {\scriptsize 0.029}}$ & $\mathbf{0.474 \pm {\scriptsize 0.012}}$ & $\mathbf{0.727 \pm {\scriptsize 0.006}}$ & $0.484 \pm {\scriptsize 0.016}$ \\
\ac{mLSTM} & 9s & $0.464 \pm {\scriptsize 0.066}$ & $0.422 \pm {\scriptsize 0.063}$ & $0.702 \pm {\scriptsize 0.029}$ & $0.459 \pm {\scriptsize 0.036}$ \\
\hline
\ac{sLSTM} & 5s & $\mathbf{0.522 \pm {\scriptsize 0.014}}$ & $\mathbf{0.471 \pm {\scriptsize 0.044}}$ & $\mathbf{0.720 \pm {\scriptsize 0.027}}$ & $\mathbf{0.513 \pm {\scriptsize 0.016}}$ \\
\ac{sLSTM} & 7s & $0.485 \pm {\scriptsize 0.055}$ & $0.466 \pm {\scriptsize 0.091}$ & $\mathbf{0.720 \pm {\scriptsize 0.047}}$ & $\mathbf{0.493 \pm {\scriptsize 0.042}}$ \\
\ac{sLSTM} & 9s & $\mathbf{0.525 \pm {\scriptsize 0.038}}$ & $\mathbf{0.479 \pm {\scriptsize 0.063}}$ & $\mathbf{0.726 \pm {\scriptsize 0.037}}$ & $\mathbf{0.486 \pm {\scriptsize 0.039}}$ \\
\hline
\ac{mLSTM}+\ac{sLSTM} & 5s & $0.500 \pm {\scriptsize 0.017}$ & $0.440 \pm {\scriptsize 0.050}$ & $0.705 \pm {\scriptsize 0.029}$ & $0.489 \pm {\scriptsize 0.013}$ \\
\ac{mLSTM}+\ac{sLSTM} & 7s & $0.505 \pm {\scriptsize 0.028}$ & $0.443 \pm {\scriptsize 0.026}$ & $0.707 \pm {\scriptsize 0.016}$ & $0.479 \pm {\scriptsize 0.020}$ \\
\ac{mLSTM}+\ac{sLSTM} & 9s & $0.501 \pm {\scriptsize 0.011}$ & $0.431 \pm {\scriptsize 0.045}$ & $0.701 \pm {\scriptsize 0.027}$ & $0.472 \pm {\scriptsize 0.023}$ \\
\hline
\end{tabular}
\setlength{\tabcolsep}{6pt}
\end{table}

\paragraph{5\,s windows.} \ac{sLSTM} attains the best mean balanced accuracy ($52.15\% \pm 1.37\%$) with very low variability, substantially more stable than the Linear Transformer ($52.06\% \pm 4.06\%$). The combined model is stable ($50.03\% \pm 1.72\%$) but trails pure \ac{sLSTM}. \ac{mLSTM} is intermediate ($50.93\% \pm 3.14\%$).

\paragraph{7\,s windows.} \ac{mLSTM} performs best ($50.66\% \pm 2.87\%$) and is the most stable in this regime. The Linear Transformer remains competitive ($50.08\% \pm 4.29\%$), while \ac{sLSTM} dips ($48.45\% \pm 5.50\%$). The combined model is again stable and near the middle ($50.46\% \pm 2.82\%$).

\paragraph{9\,s windows.} \ac{sLSTM} rebounds to the top ($52.46\% \pm 3.81\%$); the combined model is remarkably stable ($50.10\% \pm 1.14\%$). The Linear Transformer decreases monotonically with context length ($48.45\% \pm 3.73\%$). \ac{mLSTM} degrades substantially ($46.40\% \pm 6.59\%$).

\subsection{Architectural Comparison (Across Windows)}

Figure \ref{fig:architecture_window_aucpr_macro} summarizes macro-\ac{AUCPR} performance across architectures and temporal window lengths, aggregated over five random seeds per configuration. Higher \ac{AUCPR} values indicate improved performance on rare positive classes, which is particularly relevant for clinically important but infrequent EEG events.

Across all window lengths, \ac{xLSTM}-based architectures generally outperform the Linear Transformer baseline in terms of median \ac{AUCPR}, although the magnitude and consistency of improvement vary with temporal context. At 5 s, the \ac{sLSTM} exhibits the highest median \ac{AUCPR} with comparatively limited overlap with the Linear Transformer distribution, indicating a more consistent advantage at this shorter context length. This observation aligns with the statistically significant improvement reported in Table \ref{tab:statistical_significance_5s}.

At 7 s, performance differences between architectures become less pronounced. Both \ac{mLSTM} and \ac{sLSTM} maintain competitive median \ac{AUCPR} values, but their interquartile ranges overlap substantially with the Linear Transformer, suggesting that improvements are less robust relative to between-seed variability. At 9 s, the \ac{sLSTM} and combined \ac{mLSTM}+\ac{sLSTM} models retain moderately higher medians than the baseline but again with noticeable overlap, indicating that extended temporal context does not consistently yield performance gains for any architecture in \ac{AUCPR} terms.

The combined \ac{mLSTM}+\ac{sLSTM} configuration generally shows relatively narrow performance dispersion across window lengths, indicating stable optimization behavior, although it does not achieve the highest median performance at any evaluated window. In contrast, the Linear Transformer displays larger variability, particularly at shorter windows, suggesting less stable training under the same conditions.

Taken together, these results indicate that \ac{xLSTM} architectures can provide \ac{AUCPR} benefits over Linear Transformer baselines, especially at shorter window durations. However, the extent of improvement is window-dependent, and overlapping distributions at 7 s and 9 s imply that observed gains at longer windows should be interpreted with caution.\begin{figure}[htbp]
\centering
\includegraphics[width=0.95\textwidth]{images/architecture_window_aucpr_macro_box.png}
\caption{\ac{AUCPR} Macro across architectures and temporal window configurations (5\,s, 7\,s, 9\,s). Boxplots show the distribution over five random seeds; grouped boxes correspond to window length, and colors denote the architecture. At $5\,\text{s}$, \ac{sLSTM} shows a higher median than the Linear Transformer with less overlap.}
\label{fig:architecture_window_aucpr_macro}
\end{figure}

\section{Statistical Significance Testing}

Pairwise statistical comparisons were conducted between the Linear Transformer baseline and the best-performing \ac{xLSTM} variant at each window length. As summarized in Section~\ref{sec:statistical_testing}, both paired t-tests and Wilcoxon signed-rank tests were applied across seeds ($n=5$, $\alpha=0.05$). At 5\,s and 9\,s windows, \ac{sLSTM} was the top-performing variant; at 7\,s, comparisons were made against \ac{mLSTM}. Tests were computed on per-seed results; tables report mean values for readability. Given the limited sample size, p-values should be interpreted cautiously and alongside effect magnitudes.

\subsection{5-Second Windows}

\begin{table}[htbp]
\centering
\caption{Linear Transformer vs.\ \ac{sLSTM} at 5\,s (balanced-accuracy checkpoints). Positive differences favor \ac{sLSTM}.}
\label{tab:statistical_significance_5s}
\footnotesize
\begin{tabular}{l | c | c | c | c | c}
\textbf{Metric} & \textbf{Lin. Transf.} & \textbf{\ac{sLSTM}} & \textbf{t-test p} & \textbf{Wilcoxon p} & \textbf{Interpretation} \\
\hline
Balanced Accuracy  & 0.521 & 0.522 & 0.974 & 0.812 & Equivalent \\
Cohen's $\kappa$    & 0.457 & 0.471 & 0.583 & 0.438 & Higher, not significant \\
Weighted F1         & 0.712 & 0.720 & 0.651 & 0.438 & Higher, not significant \\
Accuracy            & 0.689 & 0.711 & 0.252 & 0.312 & Higher, not significant \\
AUCPR (Macro)       & 0.468 & 0.513 & \textbf{0.005} & 0.062 & \textbf{Significant (t-test)} \\
\hline
\end{tabular}
\end{table}

\noindent\textbf{Summary.} At 5\,s, \ac{sLSTM} provides a statistically significant improvement in AUCPR, indicating enhanced ranking performance for rare classes. Other metrics remain statistically equivalent.

\subsection{7-Second Windows}

\begin{table}[htbp]
\centering
\caption{Linear Transformer vs.\ \ac{mLSTM} at 7\,s (balanced-accuracy checkpoints). Positive differences favor \ac{mLSTM}.}
\label{tab:statistical_significance_7s}
\footnotesize
\begin{tabular}{l | c | c | c | c | c}
\textbf{Metric} & \textbf{Lin. Transf.} & \textbf{\ac{mLSTM}} & \textbf{t-test p} & \textbf{Wilcoxon p} & \textbf{Interpretation} \\
\hline
Balanced Accuracy  & 0.501 & 0.507 & 0.749 & 0.812 & Equivalent \\
Cohen's $\kappa$    & 0.465 & 0.474 & 0.394 & 0.438 & Equivalent \\
Weighted F1         & 0.717 & 0.727 & 0.298 & 0.312 & Equivalent \\
Accuracy            & 0.689 & 0.711 & 0.215 & 0.188 & Higher, not significant \\
AUCPR (Macro)       & 0.465 & 0.484 & 0.315 & 0.438 & Higher, not significant \\
\hline
\end{tabular}
\end{table}

\noindent\textbf{Summary.} At 7\,s, \ac{mLSTM} achieves slightly higher mean performance across most metrics, but differences relative to the Linear Transformer remain statistically non-significant.

\subsection{9-Second Windows}

\begin{table}[htbp]
\centering
\caption{Linear Transformer vs.\ \ac{sLSTM} at 9\,s (balanced-accuracy checkpoints). Positive differences favor \ac{sLSTM}.}
\label{tab:statistical_significance_9s}
\footnotesize
\begin{tabular}{l | c | c | c | c | c}
\textbf{Metric} & \textbf{Lin. Transf.} & \textbf{\ac{sLSTM}} & \textbf{t-test p} & \textbf{Wilcoxon p} & \textbf{Interpretation} \\
\hline
Balanced Accuracy  & 0.485 & 0.525 & 0.119 & 0.188 & Higher, not significant \\
Cohen's $\kappa$    & 0.425 & 0.479 & 0.055 & 0.062 & Marginal \\
Weighted F1         & 0.691 & 0.726 & 0.073 & 0.062 & Marginal \\
Accuracy            & 0.667 & 0.714 & \textbf{0.029} & 0.062 & \textbf{Significant (t-test)} \\
AUCPR (Macro)       & 0.453 & 0.486 & 0.115 & 0.062 & Higher, not significant \\
\hline
\end{tabular}
\end{table}

\noindent\textbf{Summary.} At 9\,s, \ac{sLSTM} yields a statistically significant improvement in overall accuracy and consistent, though marginal, gains in $\kappa$ and weighted F1. Balanced accuracy improves in magnitude but does not reach significance.

\subsection{Interpretation}

\textbf{Window-dependent effects.} Performance benefits are window-specific. \ac{sLSTM} provides significant AUCPR improvements at short windows (5\,s) and significant accuracy gains at long windows (9\,s), whereas \ac{mLSTM} performs best numerically at 7\,s without reaching statistical significance.

\textbf{Metric-specific effects.} Architectural changes primarily influence ranking quality and confidence calibration rather than class-balanced decision boundaries, as reflected by improvements in AUCPR and accuracy rather than balanced accuracy.

\textbf{Practical relevance.} Effect sizes remain modest (generally $<5\%$ absolute). Combined with increased computational cost of \ac{xLSTM} variants, practical deployment benefits appear constrained unless architectural refinements further stabilize and amplify performance gains.

\section{Computational Efficiency Analysis}
\label{sec:efficiency}

Training time is a key practical constraint. Figure~\ref{fig:architecture_window_runtime_bar} visualizes the mean wall-clock time per complete run (\textbf{100 epochs}) under identical hardware and batch sizes for easier comparison.

\begin{figure}[htbp]
\centering
\includegraphics[width=0.95\textwidth]{images/architecture_window_runtime_bar_hours.png}
\caption{Training runtime across architectures and temporal window configurations. Bars show mean runtime per complete run (100 epochs) over five seeds; error bars indicate the standard deviation (in minutes). The y-axis is shown in hours.}
\label{fig:architecture_window_runtime_bar}
\end{figure}

\paragraph{Linear Transformer.} Fastest across all settings (1h49m–3h15m), delivering 2.3–3.4\(\times\) speedups over \ac{xLSTM}. Very low variance indicates predictable compute. Runtime scales near-linearly with window length.

\paragraph{\ac{sLSTM}.} Most efficient \ac{xLSTM} (4h10m–6h40m) and the best performance–efficiency compromise. Variance is moderate; scaling with window length is sublinear.

\paragraph{\ac{mLSTM}.} Most expensive (5h30m–10h57m) with superlinear scaling and high variance at 9\,s, implying less predictable requirements as context grows.

\paragraph{Combined.} Nearly identical to \ac{mLSTM} in cost, confirming that the 7:1 mixture is dominated by \ac{mLSTM} blocks. Variance is somewhat more controlled than pure \ac{mLSTM} at 9\,s.

\section{Class-Specific Performance Analysis}

While aggregate metrics provide overall performance assessment, understanding class-level behavior is essential for evaluating model utility in clinical settings where minority classes often represent critical pathological events. We analyze normalized confusion matrices aggregated across all five random seeds, where diagonal entries represent per-class recall (sensitivity) and off-diagonal entries reveal systematic misclassification patterns. Values are reported as mean~\(\pm\)~standard deviation to capture both typical performance and cross-seed stability.

\subsection{Per-Class Recall Patterns}

\subsubsection{Baseline Performance at 5-Second Windows}

\begin{table}[htbp]
\centering
\caption{Normalized confusion matrix for Linear Transformer at 5\,s windows, averaged over five seeds. Diagonal entries (bold) show per-class recall; off-diagonal entries show misclassification patterns.}
\label{tab:cm_linear_5s}
\small
\setlength{\tabcolsep}{2pt}
\begin{tabular}{lcccccc}
\toprule
\textbf{True\textbackslash{}Pred} & \textbf{Class 0} & \textbf{Class 1} & \textbf{Class 2} & \textbf{Class 3} & \textbf{Class 4} & \textbf{Class 5} \\
\midrule
\textbf{Class 0} & \textbf{0.082}$\pm{\scriptsize 0.164}$ & 0.259$\pm{\scriptsize 0.135}$ & 0.140$\pm{\scriptsize 0.109}$ & 0.070$\pm{\scriptsize 0.067}$ & 0.123$\pm{\scriptsize 0.150}$ & 0.325$\pm{\scriptsize 0.082}$ \\
\textbf{Class 1} & 0.027$\pm{\scriptsize 0.027}$ & \textbf{0.411}$\pm{\scriptsize 0.119}$ & 0.201$\pm{\scriptsize 0.035}$ & 0.010$\pm{\scriptsize 0.009}$ & 0.229$\pm{\scriptsize 0.114}$ & 0.122$\pm{\scriptsize 0.061}$ \\
\textbf{Class 2} & 0.047$\pm{\scriptsize 0.022}$ & 0.144$\pm{\scriptsize 0.038}$ & \textbf{0.567}$\pm{\scriptsize 0.064}$ & 0.065$\pm{\scriptsize 0.041}$ & 0.063$\pm{\scriptsize 0.035}$ & 0.113$\pm{\scriptsize 0.037}$ \\
\textbf{Class 3} & 0.005$\pm{\scriptsize 0.003}$ & 0.002$\pm{\scriptsize 0.005}$ & 0.022$\pm{\scriptsize 0.019}$ & \textbf{0.664}$\pm{\scriptsize 0.046}$ & 0.086$\pm{\scriptsize 0.045}$ & 0.221$\pm{\scriptsize 0.034}$ \\
\textbf{Class 4} & 0.007$\pm{\scriptsize 0.006}$ & 0.006$\pm{\scriptsize 0.005}$ & 0.029$\pm{\scriptsize 0.029}$ & 0.038$\pm{\scriptsize 0.020}$ & \textbf{0.605}$\pm{\scriptsize 0.099}$ & 0.316$\pm{\scriptsize 0.088}$ \\
\textbf{Class 5} & 0.029$\pm{\scriptsize 0.018}$ & 0.011$\pm{\scriptsize 0.005}$ & 0.047$\pm{\scriptsize 0.011}$ & 0.029$\pm{\scriptsize 0.014}$ & 0.089$\pm{\scriptsize 0.034}$ & \textbf{0.794}$\pm{\scriptsize 0.030}$ \\
\bottomrule
\end{tabular}
\setlength{\tabcolsep}{6pt}
\end{table}

The baseline reveals severe performance stratification by class frequency. Class~0 (1.93\% of training samples) achieves only 8.2\% mean recall with standard deviation (16.4\%) exceeding the mean, indicating highly unstable predictions dominated by random seed effects. The majority of Class~0 instances are misclassified as Class~5 (32.5\%) or Class~1 (25.9\%), suggesting the model defaults to more frequent classes when confronted with rare patterns.

Moderate-frequency classes show intermediate performance: Class~1 (41.1\% recall), Class~2 (56.7\%), Class~3 (66.4\%), and Class~4 (60.5\%). Class~3 and Class~4 exhibit substantial confusion with Class~5 (22.1\% and 31.6\% respectively), reflecting the challenge of discriminating artifact patterns from background activity. The majority class (Class~5, 66.77\% of samples) achieves 79.4\% recall with low variance (3.0\%), demonstrating stable high performance on frequent patterns.

\subsubsection{xLSTM Performance at Optimal Configurations}

\begin{table}[htbp]
\centering
\caption{Normalized confusion matrix for \ac{sLSTM} at 5\,s windows, averaged over five seeds.}
\label{tab:cm_slstm_5s}
\small
\setlength{\tabcolsep}{2pt}
\begin{tabular}{lcccccc}
\toprule
\textbf{True\textbackslash{}Pred} & \textbf{Class 0} & \textbf{Class 1} & \textbf{Class 2} & \textbf{Class 3} & \textbf{Class 4} & \textbf{Class 5} \\
\midrule
\textbf{Class 0} & \textbf{0.051}$\pm{\scriptsize 0.078}$ & 0.400$\pm{\scriptsize 0.058}$ & 0.238$\pm{\scriptsize 0.022}$ & 0.056$\pm{\scriptsize 0.023}$ & 0.062$\pm{\scriptsize 0.059}$ & 0.194$\pm{\scriptsize 0.085}$ \\
\textbf{Class 1} & 0.027$\pm{\scriptsize 0.018}$ & \textbf{0.420}$\pm{\scriptsize 0.143}$ & 0.184$\pm{\scriptsize 0.015}$ & 0.007$\pm{\scriptsize 0.005}$ & 0.108$\pm{\scriptsize 0.063}$ & 0.255$\pm{\scriptsize 0.090}$ \\
\textbf{Class 2} & 0.093$\pm{\scriptsize 0.060}$ & 0.116$\pm{\scriptsize 0.062}$ & \textbf{0.525}$\pm{\scriptsize 0.095}$ & 0.047$\pm{\scriptsize 0.044}$ & 0.077$\pm{\scriptsize 0.025}$ & 0.142$\pm{\scriptsize 0.052}$ \\
\textbf{Class 3} & 0.004$\pm{\scriptsize 0.008}$ & 0.000$\pm{\scriptsize 0.000}$ & 0.012$\pm{\scriptsize 0.019}$ & \textbf{0.664}$\pm{\scriptsize 0.120}$ & 0.097$\pm{\scriptsize 0.042}$ & 0.224$\pm{\scriptsize 0.086}$ \\
\textbf{Class 4} & 0.001$\pm{\scriptsize 0.001}$ & 0.010$\pm{\scriptsize 0.007}$ & 0.011$\pm{\scriptsize 0.006}$ & 0.030$\pm{\scriptsize 0.016}$ & \textbf{0.643}$\pm{\scriptsize 0.057}$ & 0.305$\pm{\scriptsize 0.053}$ \\
\textbf{Class 5} & 0.017$\pm{\scriptsize 0.018}$ & 0.018$\pm{\scriptsize 0.007}$ & 0.030$\pm{\scriptsize 0.013}$ & 0.011$\pm{\scriptsize 0.006}$ & 0.097$\pm{\scriptsize 0.020}$ & \textbf{0.826}$\pm{\scriptsize 0.015}$ \\
\bottomrule
\end{tabular}
\setlength{\tabcolsep}{6pt}
\end{table}

Comparing \ac{sLSTM} to the baseline reveals minimal architectural impact on class-specific discrimination. Class~0 recall actually decreases (5.1\% vs.\ 8.2\%), though neither value represents clinically useful performance. The primary difference emerges in the majority class: \ac{sLSTM} achieves 82.6\% recall on Class~5 with substantially lower variance (1.5\% vs.\ 3.0\%), consistent with its superior overall optimization stability. Intermediate classes show nearly identical recall patterns, with differences well within cross-seed variance.

\subsubsection{Performance at 7-Second Windows}

At 7\,s windows, architectural differences become more pronounced. Table~\ref{tab:cm_linear_7s} shows the Linear Transformer baseline at this window length.

\begin{table}[htbp]
\centering
\caption{Normalized confusion matrix for Linear Transformer at 7\,s windows, averaged over five seeds.}
\label{tab:cm_linear_7s}
\small
\setlength{\tabcolsep}{2pt}
\begin{tabular}{lcccccc}
\toprule
\textbf{True\textbackslash{}Pred} & \textbf{Class 0} & \textbf{Class 1} & \textbf{Class 2} & \textbf{Class 3} & \textbf{Class 4} & \textbf{Class 5} \\
\midrule
\textbf{Class 0} & \textbf{0.066}$\pm{\scriptsize 0.089}$ & 0.383$\pm{\scriptsize 0.130}$ & 0.086$\pm{\scriptsize 0.039}$ & 0.090$\pm{\scriptsize 0.043}$ & 0.141$\pm{\scriptsize 0.155}$ & 0.234$\pm{\scriptsize 0.036}$ \\
\textbf{Class 1} & 0.028$\pm{\scriptsize 0.028}$ & \textbf{0.485}$\pm{\scriptsize 0.117}$ & 0.191$\pm{\scriptsize 0.070}$ & 0.008$\pm{\scriptsize 0.009}$ & 0.235$\pm{\scriptsize 0.114}$ & 0.053$\pm{\scriptsize 0.017}$ \\
\textbf{Class 2} & 0.079$\pm{\scriptsize 0.071}$ & 0.181$\pm{\scriptsize 0.096}$ & \textbf{0.466}$\pm{\scriptsize 0.100}$ & 0.088$\pm{\scriptsize 0.049}$ & 0.065$\pm{\scriptsize 0.016}$ & 0.119$\pm{\scriptsize 0.061}$ \\
\textbf{Class 3} & 0.002$\pm{\scriptsize 0.004}$ & 0.009$\pm{\scriptsize 0.009}$ & 0.016$\pm{\scriptsize 0.013}$ & \textbf{0.593}$\pm{\scriptsize 0.082}$ & 0.104$\pm{\scriptsize 0.051}$ & 0.275$\pm{\scriptsize 0.061}$ \\
\textbf{Class 4} & 0.002$\pm{\scriptsize 0.003}$ & 0.011$\pm{\scriptsize 0.009}$ & 0.013$\pm{\scriptsize 0.011}$ & 0.071$\pm{\scriptsize 0.028}$ & \textbf{0.606}$\pm{\scriptsize 0.073}$ & 0.296$\pm{\scriptsize 0.092}$ \\
\textbf{Class 5} & 0.022$\pm{\scriptsize 0.009}$ & 0.015$\pm{\scriptsize 0.014}$ & 0.034$\pm{\scriptsize 0.010}$ & 0.033$\pm{\scriptsize 0.017}$ & 0.107$\pm{\scriptsize 0.037}$ & \textbf{0.788}$\pm{\scriptsize 0.066}$ \\
\bottomrule
\end{tabular}
\setlength{\tabcolsep}{6pt}
\end{table}

The Linear Transformer at 7\,s maintains relatively stable performance across most classes, with Class~1 achieving 48.5\% recall, slightly higher than at 5\,s. However, Class~0 performance remains poor at 6.6\%, with high variance indicating continued instability on rare classes.

\begin{table}[htbp]
\centering
\caption{Normalized confusion matrix for \ac{mLSTM} at 7\,s windows, averaged over five seeds. Best overall performance at this window length.}
\label{tab:cm_mlstm_7s}
\small
\setlength{\tabcolsep}{2pt}
\begin{tabular}{lcccccc}
\toprule
\textbf{True\textbackslash{}Pred} & \textbf{Class 0} & \textbf{Class 1} & \textbf{Class 2} & \textbf{Class 3} & \textbf{Class 4} & \textbf{Class 5} \\
\midrule
\textbf{Class 0} & \textbf{0.083}$\pm{\scriptsize 0.108}$ & 0.425$\pm{\scriptsize 0.115}$ & 0.044$\pm{\scriptsize 0.031}$ & 0.139$\pm{\scriptsize 0.061}$ & 0.019$\pm{\scriptsize 0.022}$ & 0.289$\pm{\scriptsize 0.053}$ \\
\textbf{Class 1} & 0.030$\pm{\scriptsize 0.028}$ & \textbf{0.532}$\pm{\scriptsize 0.091}$ & 0.141$\pm{\scriptsize 0.037}$ & 0.010$\pm{\scriptsize 0.006}$ & 0.114$\pm{\scriptsize 0.078}$ & 0.171$\pm{\scriptsize 0.112}$ \\
\textbf{Class 2} & 0.076$\pm{\scriptsize 0.053}$ & 0.137$\pm{\scriptsize 0.030}$ & \textbf{0.483}$\pm{\scriptsize 0.087}$ & 0.053$\pm{\scriptsize 0.028}$ & 0.100$\pm{\scriptsize 0.026}$ & 0.150$\pm{\scriptsize 0.043}$ \\
\textbf{Class 3} & 0.002$\pm{\scriptsize 0.005}$ & 0.007$\pm{\scriptsize 0.009}$ & 0.016$\pm{\scriptsize 0.016}$ & \textbf{0.579}$\pm{\scriptsize 0.036}$ & 0.160$\pm{\scriptsize 0.021}$ & 0.235$\pm{\scriptsize 0.052}$ \\
\textbf{Class 4} & 0.005$\pm{\scriptsize 0.005}$ & 0.011$\pm{\scriptsize 0.010}$ & 0.012$\pm{\scriptsize 0.012}$ & 0.008$\pm{\scriptsize 0.009}$ & \textbf{0.546}$\pm{\scriptsize 0.077}$ & 0.417$\pm{\scriptsize 0.064}$ \\
\textbf{Class 5} & 0.028$\pm{\scriptsize 0.017}$ & 0.024$\pm{\scriptsize 0.013}$ & 0.020$\pm{\scriptsize 0.007}$ & 0.006$\pm{\scriptsize 0.003}$ & 0.106$\pm{\scriptsize 0.010}$ & \textbf{0.815}$\pm{\scriptsize 0.020}$ \\
\bottomrule
\end{tabular}
\setlength{\tabcolsep}{6pt}
\end{table}

\ac{mLSTM} at 7\,s achieves the highest Class~1 recall across all configurations (53.2\%), suggesting matrix memory better captures patterns in this moderately frequent class at intermediate sequence lengths. However, Class~0 remains problematic (8.3\%), and Class~4 recall drops to 54.6\% with increased confusion with Class~5 (41.7\%), the highest misclassification rate observed for this class pair. This trade-off pattern—improved performance on some classes at the expense of others—explains why \ac{mLSTM}'s aggregate balanced accuracy remains competitive without clearly dominating.

\subsubsection{Performance at 9-Second Windows}

At 9\,s windows, the performance landscape shifts again. Table~\ref{tab:cm_linear_9s} shows the Linear Transformer baseline at this extended window length.

\begin{table}[htbp]
\centering
\caption{Normalized confusion matrix for Linear Transformer at 9\,s windows, averaged over five seeds.}
\label{tab:cm_linear_9s}
\small
\setlength{\tabcolsep}{2pt}
\begin{tabular}{lcccccc}
\toprule
\textbf{True\textbackslash{}Pred} & \textbf{Class 0} & \textbf{Class 1} & \textbf{Class 2} & \textbf{Class 3} & \textbf{Class 4} & \textbf{Class 5} \\
\midrule
\textbf{Class 0} & \textbf{0.017}$\pm{\scriptsize 0.028}$ & 0.349$\pm{\scriptsize 0.121}$ & 0.240$\pm{\scriptsize 0.069}$ & 0.078$\pm{\scriptsize 0.061}$ & 0.031$\pm{\scriptsize 0.048}$ & 0.286$\pm{\scriptsize 0.049}$ \\
\textbf{Class 1} & 0.007$\pm{\scriptsize 0.005}$ & \textbf{0.404}$\pm{\scriptsize 0.183}$ & 0.270$\pm{\scriptsize 0.076}$ & 0.008$\pm{\scriptsize 0.008}$ & 0.175$\pm{\scriptsize 0.102}$ & 0.135$\pm{\scriptsize 0.076}$ \\
\textbf{Class 2} & 0.041$\pm{\scriptsize 0.024}$ & 0.162$\pm{\scriptsize 0.075}$ & \textbf{0.512}$\pm{\scriptsize 0.117}$ & 0.099$\pm{\scriptsize 0.060}$ & 0.053$\pm{\scriptsize 0.040}$ & 0.134$\pm{\scriptsize 0.068}$ \\
\textbf{Class 3} & 0.009$\pm{\scriptsize 0.012}$ & 0.007$\pm{\scriptsize 0.010}$ & 0.019$\pm{\scriptsize 0.010}$ & \textbf{0.616}$\pm{\scriptsize 0.153}$ & 0.068$\pm{\scriptsize 0.048}$ & 0.281$\pm{\scriptsize 0.130}$ \\
\textbf{Class 4} & 0.008$\pm{\scriptsize 0.005}$ & 0.008$\pm{\scriptsize 0.012}$ & 0.022$\pm{\scriptsize 0.012}$ & 0.061$\pm{\scriptsize 0.041}$ & \textbf{0.584}$\pm{\scriptsize 0.054}$ & 0.317$\pm{\scriptsize 0.052}$ \\
\textbf{Class 5} & 0.031$\pm{\scriptsize 0.019}$ & 0.012$\pm{\scriptsize 0.009}$ & 0.049$\pm{\scriptsize 0.018}$ & 0.036$\pm{\scriptsize 0.028}$ & 0.099$\pm{\scriptsize 0.020}$ & \textbf{0.774}$\pm{\scriptsize 0.044}$ \\
\bottomrule
\end{tabular}
\setlength{\tabcolsep}{6pt}
\end{table}

The Linear Transformer at 9\,s shows continued degradation in Class~0 (1.7\% recall) and substantial decline in Class~5 (77.4\%, down from 79.4\% at 5\,s and 78.8\% at 7\,s). Class~1 recall also drops to 40.4\%, indicating that extended context length degrades the baseline's ability to maintain performance across classes. This monotonic decline with window length suggests attention mechanisms struggle to effectively leverage longer temporal contexts in this task.

\begin{table}[htbp]
\centering
\caption{Normalized confusion matrix for \ac{sLSTM} at 9\,s windows, averaged over five seeds. Best overall balanced accuracy at this window length.}
\label{tab:cm_slstm_9s}
\small
\setlength{\tabcolsep}{2pt}
\begin{tabular}{lcccccc}
\toprule
\textbf{True\textbackslash{}Pred} & \textbf{Class 0} & \textbf{Class 1} & \textbf{Class 2} & \textbf{Class 3} & \textbf{Class 4} & \textbf{Class 5} \\
\midrule
\textbf{Class 0} & \textbf{0.142}$\pm{\scriptsize 0.094}$ & 0.371$\pm{\scriptsize 0.144}$ & 0.223$\pm{\scriptsize 0.151}$ & 0.096$\pm{\scriptsize 0.061}$ & 0.014$\pm{\scriptsize 0.022}$ & 0.153$\pm{\scriptsize 0.075}$ \\
\textbf{Class 1} & 0.047$\pm{\scriptsize 0.058}$ & \textbf{0.444}$\pm{\scriptsize 0.173}$ & 0.238$\pm{\scriptsize 0.118}$ & 0.007$\pm{\scriptsize 0.010}$ & 0.078$\pm{\scriptsize 0.041}$ & 0.185$\pm{\scriptsize 0.153}$ \\
\textbf{Class 2} & 0.047$\pm{\scriptsize 0.030}$ & 0.170$\pm{\scriptsize 0.015}$ & \textbf{0.516}$\pm{\scriptsize 0.089}$ & 0.059$\pm{\scriptsize 0.043}$ & 0.061$\pm{\scriptsize 0.041}$ & 0.146$\pm{\scriptsize 0.071}$ \\
\textbf{Class 3} & 0.020$\pm{\scriptsize 0.038}$ & 0.002$\pm{\scriptsize 0.004}$ & 0.006$\pm{\scriptsize 0.011}$ & \textbf{0.638}$\pm{\scriptsize 0.170}$ & 0.075$\pm{\scriptsize 0.080}$ & 0.258$\pm{\scriptsize 0.106}$ \\
\textbf{Class 4} & 0.011$\pm{\scriptsize 0.008}$ & 0.015$\pm{\scriptsize 0.009}$ & 0.015$\pm{\scriptsize 0.011}$ & 0.038$\pm{\scriptsize 0.039}$ & \textbf{0.576}$\pm{\scriptsize 0.055}$ & 0.345$\pm{\scriptsize 0.048}$ \\
\textbf{Class 5} & 0.022$\pm{\scriptsize 0.010}$ & 0.029$\pm{\scriptsize 0.006}$ & 0.027$\pm{\scriptsize 0.017}$ & 0.013$\pm{\scriptsize 0.014}$ & 0.078$\pm{\scriptsize 0.036}$ & \textbf{0.831}$\pm{\scriptsize 0.045}$ \\
\bottomrule
\end{tabular}
\setlength{\tabcolsep}{6pt}
\end{table}

At 9\,s windows, \ac{sLSTM} achieves 14.2\% recall on Class~0, the highest observed across all configurations. However, the high standard deviation (9.4\%, representing 66\% of the mean) indicates this improvement is driven by occasional favorable seeds rather than systematic architectural advantage. Class~5 recall remains strong (83.1\%), and intermediate classes show performance comparable to 5\,s configurations. The recovery from the 7\,s performance dip suggests \ac{sLSTM}'s scalar memory mechanisms operate more effectively at both short and extended sequence lengths, with a performance valley at intermediate scales.

\subsection{Statistical Comparison of Architectural Effects}

To rigorously assess whether observed class-level differences reflect genuine architectural advantages or random variation, we perform paired statistical tests comparing Linear Transformer and \ac{sLSTM} at each window length.

\subsubsection{5-Second Windows}

\begin{table}[htbp]
\centering
\caption{Statistical comparison of per-class recall: Linear Transformer vs.\ \ac{sLSTM} at 5\,s windows. Tests are paired across seeds; \(\alpha=0.05\).}
\label{tab:class_comparison_5s}
\small
\setlength{\tabcolsep}{3pt}
\begin{tabular}{l|c|c|c|c|c}
\hline
\textbf{Class} & \textbf{Linear} & \textbf{\ac{sLSTM}} & \textbf{t-stat} & \textbf{p (t-test)} & \textbf{p (Wilcoxon)} \\
\hline
Class 0 & 0.082$\pm{\scriptsize 0.164}$ & 0.051$\pm{\scriptsize 0.078}$ & 0.33 & 0.756 & 1.000 \\
Class 1 & 0.411$\pm{\scriptsize 0.119}$ & 0.420$\pm{\scriptsize 0.143}$ & $-0.13$ & 0.902 & 0.812 \\
Class 2 & 0.567$\pm{\scriptsize 0.064}$ & 0.525$\pm{\scriptsize 0.095}$ & 0.92 & 0.410 & 0.312 \\
Class 3 & 0.664$\pm{\scriptsize 0.046}$ & 0.664$\pm{\scriptsize 0.120}$ & 0.01 & 0.989 & 1.000 \\
Class 4 & 0.605$\pm{\scriptsize 0.099}$ & 0.643$\pm{\scriptsize 0.057}$ & $-0.71$ & 0.515 & 0.812 \\
Class 5 & 0.794$\pm{\scriptsize 0.030}$ & 0.826$\pm{\scriptsize 0.015}$ & $-2.38$ & 0.076 & 0.062 \\
\hline
\end{tabular}
\setlength{\tabcolsep}{6pt}
\end{table}

At 5\,s windows, no class-level differences reach statistical significance at \(\alpha=0.05\). Class~5 shows the strongest trend toward \ac{sLSTM} advantage (recall difference: +3.2 percentage points, p=0.076), approaching but not achieving significance. This near-significant improvement in majority-class recall, combined with marginal differences in other classes, explains why \ac{sLSTM} achieves significantly better \ac{AUCPR} Macro (ranking quality; Section~\ref{tab:statistical_significance_5s}) without significant balanced accuracy improvement. The architectural benefit manifests primarily in confidence calibration for the majority class rather than improved discrimination of minority classes.

\textbf{In summary,} at 5\,s windows, \ac{sLSTM} provides no significant class-level improvements but shows trends toward better majority-class performance and ranking quality.

\subsubsection{7-Second Windows}

\begin{table}[htbp]
\centering
\caption{Statistical comparison of per-class recall: Linear Transformer vs.\ \ac{mLSTM} at 7\,s windows.}
\label{tab:class_comparison_7s}
\small
\setlength{\tabcolsep}{3pt}
\begin{tabular}{l|c|c|c|c|c}
\hline
\textbf{Class} & \textbf{Linear} & \textbf{\ac{mLSTM}} & \textbf{t-stat} & \textbf{p (t-test)} & \textbf{p (Wilcoxon)} \\
\hline
Class 0 & 0.066$\pm{\scriptsize 0.089}$ & 0.083$\pm{\scriptsize 0.108}$ & $-0.29$ & 0.789 & 0.625 \\
Class 1 & 0.485$\pm{\scriptsize 0.117}$ & 0.532$\pm{\scriptsize 0.091}$ & $-1.05$ & 0.354 & 0.438 \\
Class 2 & 0.466$\pm{\scriptsize 0.100}$ & 0.483$\pm{\scriptsize 0.087}$ & $-0.32$ & 0.763 & 0.812 \\
Class 3 & 0.593$\pm{\scriptsize 0.082}$ & 0.579$\pm{\scriptsize 0.036}$ & 0.32 & 0.764 & 1.000 \\
Class 4 & 0.606$\pm{\scriptsize 0.073}$ & 0.546$\pm{\scriptsize 0.077}$ & 0.96 & 0.390 & 0.812 \\
Class 5 & 0.788$\pm{\scriptsize 0.066}$ & 0.815$\pm{\scriptsize 0.020}$ & $-0.80$ & 0.469 & 0.625 \\
\hline
\end{tabular}
\setlength{\tabcolsep}{6pt}
\end{table}

At 7\,s windows, no class-level differences approach significance (all p>0.35). \ac{mLSTM} shows numerically higher recall for Classes~0, 1, 2, and 5, with the largest improvement in Class~1 (+4.7 percentage points, p=0.354). However, Class~4 recall decreases substantially ($-6.0$ percentage points, p=0.390), indicating a trade-off where matrix memory improves performance on some classes at the expense of others. This mixed pattern explains why \ac{mLSTM}'s aggregate balanced accuracy remains competitive without clearly dominating (Section~\ref{tab:statistical_significance_7s}). The higher Class~1 recall (53.2\%, the best across all configurations) suggests matrix memory mechanisms are particularly effective for this moderately frequent class at intermediate sequence lengths.

\textbf{In summary,} at 7\,s windows, \ac{mLSTM} achieves the best Class~1 recall across all configurations but shows trade-offs in Class~4, resulting in no significant aggregate advantages over the Linear Transformer baseline.

\subsubsection{9-Second Windows}

\begin{table}[htbp]
\centering
\caption{Statistical comparison of per-class recall: Linear Transformer vs.\ \ac{sLSTM} at 9\,s windows.}
\label{tab:class_comparison_9s}
\small
\setlength{\tabcolsep}{3pt}
\begin{tabular}{l|c|c|c|c|c}
\hline
\textbf{Class} & \textbf{Linear} & \textbf{\ac{sLSTM}} & \textbf{t-stat} & \textbf{p (t-test)} & \textbf{p (Wilcoxon)} \\
\hline
Class 0 & 0.017$\pm{\scriptsize 0.028}$ & 0.142$\pm{\scriptsize 0.094}$ & $-2.84$ & \textbf{0.047}* & 0.125 \\
Class 1 & 0.404$\pm{\scriptsize 0.183}$ & 0.444$\pm{\scriptsize 0.173}$ & $-0.42$ & 0.698 & 0.812 \\
Class 2 & 0.512$\pm{\scriptsize 0.117}$ & 0.516$\pm{\scriptsize 0.089}$ & $-0.07$ & 0.948 & 1.000 \\
Class 3 & 0.616$\pm{\scriptsize 0.153}$ & 0.638$\pm{\scriptsize 0.170}$ & $-0.22$ & 0.839 & 1.000 \\
Class 4 & 0.584$\pm{\scriptsize 0.054}$ & 0.576$\pm{\scriptsize 0.055}$ & 0.22 & 0.834 & 1.000 \\
Class 5 & 0.774$\pm{\scriptsize 0.044}$ & 0.831$\pm{\scriptsize 0.045}$ & $-3.97$ & \textbf{0.016}* & 0.062 \\
\hline
\end{tabular}
\setlength{\tabcolsep}{6pt}
\end{table}

At 9\,s windows, \ac{sLSTM} achieves statistically significant improvements in Class~0 (p=0.047; +12.5 percentage points) and Class~5 (p=0.016; +5.7 percentage points). The Class~0 improvement is substantial in relative terms (8.3× increase from 1.7\% to 14.2\%), representing the first instance where any architecture achieves double-digit recall on the rarest class. However, the clinical utility remains questionable: 14.2\% recall means 85.8\% of minority-class instances are still misclassified. The high standard deviation (9.4\%) relative to the mean indicates this improvement is highly variable across seeds, with the Wilcoxon test (p=0.125) not reaching significance, suggesting potential sensitivity to outlier seeds.

The Class~5 improvement is more robust, with both parametric and non-parametric tests approaching or reaching significance. This majority-class advantage combines with marginal improvements across intermediate classes (Classes~1–3 show small non-significant increases) to produce the significant overall accuracy improvement observed in aggregate metrics (Section~\ref{tab:statistical_significance_9s}). Notably, Class~4 shows no improvement, maintaining the pattern where extended context benefits some but not all classes.

\textbf{In summary,} at 9\,s windows, \ac{sLSTM} achieves significant improvements in the rarest class (Class~0) and the majority class (Class~5), driving its overall performance recovery at extended context lengths.

\section{Summary of Key Findings}

\subsection{Performance Characteristics}

\begin{itemize}
\item \textbf{Best overall (balanced acc.).} \ac{sLSTM} leads at 5\,s ($52.15\% \pm 1.37\%$) and 9\,s ($52.46\% \pm 3.81\%$), but drops at 7\,s ($48.45\% \pm 5.50\%$), indicating sequence-length sensitivity.
\item \textbf{Most consistent.} The combined \ac{mLSTM}+\ac{sLSTM} is highly stable across windows ($50.03$–$50.46\%$), with the lowest variability at 9\,s ($\pm 1.14\%$), but never achieves the top score.
\item \textbf{Window-specific peak.} \ac{mLSTM} peaks at 7\,s ($50.66\% \pm 2.87\%$) and collapses at 9\,s ($46.40\% \pm 6.59\%$).
\item \textbf{Baseline competitiveness.} The Linear Transformer is competitive at 5–7\,s while training 2.3–3.4\(\times\) faster; it declines with longer context.
\end{itemize}

\subsection{Statistical Significance}

\begin{itemize}
\item \textbf{Window dependence.} \ac{sLSTM} achieves significant \ac{AUCPR} gains at 5\,s (p=0.005) and significant accuracy gains at 9\,s (p=0.029); no significant differences at 7\,s.
\item \textbf{Metric dependence.} Gains shift from ranking quality (short windows) to decision accuracy (long windows). Balanced accuracy never reaches significance.
\item \textbf{Effect sizes.} All significant improvements are modest (<5\% absolute), tempering practical impact.
\end{itemize}

\subsection{Architectural Insights}

\begin{itemize}
\item \textbf{Sequence-length sensitivity.} \ac{sLSTM} and \ac{mLSTM} prefer different operating ranges (U-shaped vs.\ inverted-U), pointing to distinct memory–timescale interactions.
\item \textbf{Hybrid trade-offs.} A fixed 7:1 \ac{mLSTM}:\ac{sLSTM} mix prioritizes stability but mutes peak performance; alternative ratios or adaptive mixing merit exploration.
\item \textbf{Compute scaling.} \ac{xLSTM} variants incur 2–3\(\times\) training cost over the baseline; \ac{mLSTM} scales superlinearly with window length.
\end{itemize}

\subsection{Practical Implications}

\begin{itemize}
\item \textbf{Deployment choice.} For this task, \ac{sLSTM} (5\,s) offers the best accuracy–stability point, whereas the Linear Transformer (5\,s) provides the fastest, most predictable training with competitive accuracy.
\item \textbf{Data imbalance dominates.} Architectural changes and longer context alone do not fix minority-class recall; future work should prioritize class-imbalance remedies (e.g., cost-sensitive learning, reweighting, resampling, augmentation).
\end{itemize}
